\documentclass[12pt]{article}
\usepackage[utf8]{inputenc}
\usepackage[T1]{fontenc}
\usepackage{hyperref}
\title{Twendi Language Resources for K-12 Curriculum Development}
\author{Compiled by Claude Sonnet 5 (claude-sonnet-5) for Andr\'as Kornai}
\date{}
\begin{document}
\maketitle

\section{Introduction}
Twendi (ISO 639-3 twn, Glottolog twen1242), also called Cambap, is a Mambiloid language of Cameroon. Wikipedia describes it as nearly extinct. It is spoken in the villages of Cambap and Sanga on the Tikar Plain by no more than about 30 people, the youngest born in the 1940s. Speakers have largely shifted to the closely related Kwanja. A 1996 survey cited by the Endangered Languages Project counted about 35 active speakers.

\section{Primary Educational Resources}
None found. Search results show no primers, textbooks or school programmes.

\section{Linguistic and Reference Materials}
\begin{itemize}
\item Wikipedia article on Twendi: \url{https://en.wikipedia.org/wiki/Twendi_language}
\item Numeral system data (Bruce Connell), Max Planck EVA numerals database: \url{https://lingweb.eva.mpg.de/channumerals/Twendi.htm}
\item Wiktionary category for Twendi (I did not check how many entries it holds): \url{https://en.wiktionary.org/wiki/Category:Twendi_language}
\end{itemize}

\section{Religious and Community Resources}
Joshua Project has a profile of the Twendi people. Its population figures (about 3,100 to 3,200) concern the ethnic group, not the language, and differ from the linguistic sources. \url{https://joshuaproject.net/people_groups/21360}

\section{Digital and Audio Resources}
Endangered Languages Project entry: \url{https://endangeredlanguages.com/lang/532}. I found no audio recordings in search results.

\section{Curriculum Development Context}
The language is moribund, and its last speakers are elderly. A school curriculum is not realistic as such. The practical priority is documentation (recordings, a wordlist, texts) with the remaining speakers, and then a modest heritage component. Kwanja is the closely related language into which the community has shifted. The resources found here are descriptive only.
\end{document}